\chapter{Présentation des résultats}
\label{chap:resultats}

\section{Scénario de bout en bout}

Le scénario suivant illustre le fonctionnement complet de la chaîne mise en
œuvre~:

\begin{enumerate}
  \item Un agent Wazuh détecte un événement suspect sur un hôte surveillé
        (par exemple, une création de compte utilisateur inhabituelle) et
        génère une alerte, indexée dans le Wazuh Indexer.
  \item Le \textit{Poller Service}, lors de son cycle suivant, récupère
        cette alerte et la soumet au \textit{Policy Engine}, qui juge sa
        sévérité suffisante pour ouvrir une investigation.
  \item L'\textit{Investigation Manager} crée l'investigation, l'associe à
        l'actif correspondant, et la transmet à l'\textit{AI Investigator}.
  \item L'agent récupère son contexte de référence (fiche de l'actif,
        alertes récentes du même agent), puis décide, par exemple,
        d'interroger \texttt{search\_alerts} pour vérifier si des alertes
        similaires ont déjà été observées sur ce même hôte, et
        \texttt{check\_ip\_reputation} si une adresse IP externe est
        impliquée.
  \item Sur la base des preuves réunies, l'agent produit un verdict
        (\texttt{TRUE\_POSITIVE}, \texttt{FALSE\_POSITIVE} ou
        \texttt{NEEDS\_REVIEW}) accompagné d'un niveau de confiance et
        d'une explication.
  \item L'analyste consulte ce verdict, la chronologie complète des actions
        effectuées par l'IA et les preuves associées depuis l'interface
        web, et peut poser des questions complémentaires à l'assistant
        conversationnel (par exemple~: \og combien d'alertes ce mois-ci
        proviennent de cet hôte~?\fg{}, désormais résolu de manière exacte
        par l'outil \texttt{count\_alerts} plutôt qu'estimé à partir de
        résultats de recherche plafonnés).
\end{enumerate}

\section{Validations fonctionnelles effectuées}

Le projet ne dispose pas, à ce stade, d'une suite de tests automatisés
(voir \cref{sec:limites}). La validation s'est appuyée sur des vérifications
fonctionnelles ciblées, effectuées à chaque étape de développement contre
les services réels (Wazuh Indexer, base MySQL, API de LLM), plutôt que sur
des données simulées. Le \cref{tab:validations} en résume les principales.

\begin{table}[H]
\centering
\caption{Principales validations fonctionnelles effectuées}
\label{tab:validations}
\begin{tabular}{@{}p{5cm}p{8cm}@{}}
\toprule
\textbf{Fonctionnalité validée} & \textbf{Méthode et résultat} \\
\midrule
Boucle d'investigation agentique & Exécution de bout en bout sur des
alertes réelles jusqu'à l'obtention d'un verdict enregistré \\
Détection des appels d'outils dupliqués & Vérifié qu'un second appel
identique réutilise le résultat mis en cache au lieu de solliciter à
nouveau Wazuh/MySQL \\
Normalisation des arguments d'outils & Reproduction du cas où le modèle
transmet un entier sous forme de chaîne (\texttt{"limit": "50"}), confirmée
corrigée après implémentation \\
Comptage exact des alertes/investigations & Comparaison entre le résultat
de \texttt{count\_alerts}/\texttt{count\_investigations} et une requête de
comptage directe~: résultats identiques \\
Authentification & Scénarios~: identifiants valides, identifiants
invalides, jeton absent, jeton expiré/invalide -- chacun renvoie le code
HTTP attendu (200 ou 401) \\
Délégation au Wazuh Indexer & Connexion réussie avec un compte Wazuh
existant, rejet d'un mot de passe erroné \\
Restriction d'accès aux données d'identité & Confirmation qu'aucun outil
disponible pour l'IA ne référence la table \texttt{users} \\
Réputation d'IP (AbuseIPDB) & Appel réel sur une IP publique connue et sur
une IP privée (court-circuitée localement, sans appel réseau) \\
Interface web & Compilation de production du frontend sans erreur,
navigation entre les écrans principaux \\
\bottomrule
\end{tabular}
\end{table}

\section{Limites constatées}
\label{sec:limites}

\begin{itemize}
  \item \textbf{Absence de suite de tests automatisés}~: les validations
        décrites ci-dessus ont été effectuées manuellement à chaque
        évolution significative ; l'ajout de tests automatisés
        (unitaires sur les outils et la boucle agentique, tests
        d'intégration sur l'API) constitue une priorité pour la suite du
        projet.
  \item \textbf{Un seul niveau d'accès}~: tout utilisateur authentifié
        dispose actuellement des mêmes droits sur l'ensemble des
        fonctionnalités ; la table \texttt{users} porte un champ de rôle
        exploitable, mais aucune vérification de permission n'est encore
        appliquée dessus.
  \item \textbf{Pas d'annulation d'une investigation en cours}~: une fois
        démarrée, une investigation s'exécute jusqu'à son terme, une
        erreur, ou un redémarrage du serveur.
  \item \textbf{Variabilité liée au modèle de langage employé}~: la
        qualité et la stabilité des verdicts dépendent du modèle choisi ;
        les modèles de plus petite taille peuvent nécessiter davantage
        d'itérations avant de converger, ce que les garde-fous
        (\cref{chap:implementation}) atténuent sans l'éliminer totalement.
  \item \textbf{Portée volontairement limitée des sources de données}~:
        l'intégration d'autres sources (Suricata, journaux \textit{cloud}
        de type CloudTrail) et de mécanismes de \textit{retrieval-augmented
        generation} (RAG) sur une base de connaissances de menaces n'a pas
        été réalisée dans cette version, l'architecture par outils
        (\texttt{Tool}) étant conçue pour rendre ces ajouts additifs.
\end{itemize}

\section{Discussion}

L'approche agentique retenue présente un avantage déterminant par rapport à
une automatisation à base de règles fixes~: elle s'adapte à des situations
non anticipées lors de la conception, en mobilisant les outils pertinents au
cas par cas plutôt qu'en suivant un scénario figé. Cette flexibilité a
toutefois un coût~: la latence et le coût monétaire d'une investigation
dépendent directement du nombre d'appels au modèle de langage, ce qui a
directement motivé plusieurs des garde-fous mis en œuvre (budget d'étapes,
détection de doublons, bornage de la taille du contexte). L'expérience du
projet montre par ailleurs qu'un fournisseur de LLM externe introduit des
contraintes opérationnelles propres (quotas de débit, formats de réponse
parfois non strictement conformes au schéma attendu) qui ont dû être
absorbées par la couche applicative plutôt que supposées absentes — la
normalisation des arguments d'outils (\cref{lst:coerce}) en est un exemple
direct. C'est également ce qui justifie la conception de l'interface
\texttt{LLMProvider} comme une abstraction remplaçable~: le projet reste en
mesure de migrer vers un modèle auto-hébergé sans remise en cause de la
logique métier.
